% GENERATED FILE. Edit canonical lessons, then run npm run generate:books.
% canonical-source-hash: fnv1a64:b0b4c66ed3716f97
% canonical-lessons: FR-C92-doux, FR-C92-sale, FR-C92-jeune, FR-C92-gentil, FR-C92-etroit

\chapter{Soft, Dirty, Young, Kind, Narrow}
\label{ch:fr-soft-dirty-young-kind-narrow}
\hlchaptermodality{92}

\begin{chapteropening}
\textbf{By the end of this chapter:} \emph{I can name soft, dirty, young, kind, narrow.}

Complete the last lesson of chapter 92: I can name soft, dirty, young, kind, narrow.
\end{chapteropening}

\section[doux]{doux — soft; mild; sweet}

\label{lesson:FR-C92-doux}



\begin{quote}
Before the new one: say the French for poor, then the French for grey.
\end{quote}

\subsection*{You'll want to know: doux}
\textbf{doux} — \textquotedblleft{}soft; mild; sweet\textquotedblright{}.

\begin{cousinweb}
\textbf{doux} $\leftarrow$ Old French \emph{dous} $\leftarrow$ Latin \emph{dulcis}, \textquotedblleft{}sweet\textquotedblright{} — inherited. (EtymDB 2.1)
\end{cousinweb}

\begin{grammarlens}[title={What you've built}]
One more word for this chapter.
\end{grammarlens}

\subsection*{Your turn}
\begin{itemize}
\raggedright
  \item \emph{Say it:} \emph{doux}
  \item \emph{Say it:} \emph{doux}, once more
  \item \emph{Say it:} say \emph{doux}
  \item \emph{Recall:} say the French for poor, then the French for grey, then say \emph{doux} again
\end{itemize}

\begin{tcolorbox}[breakable,colback=teal!4,colframe=teal!35!black,title={Before you move on}]
What does doux mean? (\textquotedblleft{}Soft; mild; sweet\textquotedblright{}.) Say it once more.
\end{tcolorbox}
\section[sale]{sale — dirty}

\label{lesson:FR-C92-sale}



\begin{quote}
Before the new one: say the French for grey, then the French for soft; mild; sweet.
\end{quote}

\subsection*{You'll want to know: sale}
\textbf{sale} — \textquotedblleft{}dirty\textquotedblright{}.

\begin{cousinweb}
\textbf{sale} $\leftarrow$ Frankish \emph{*salo}, \textquotedblleft{}dull, dirty grey\textquotedblright{} — borrowed from the Franks. (EtymDB 2.1)
\end{cousinweb}

\begin{grammarlens}[title={What you've built}]
One more word for this chapter.
\end{grammarlens}

\subsection*{Your turn}
\begin{itemize}
\raggedright
  \item \emph{Say it:} \emph{sale}
  \item \emph{Say it:} \emph{sale}, once more
  \item \emph{Say it:} say \emph{sale}
  \item \emph{Recall:} say the French for grey, then the French for soft; mild; sweet, then say \emph{sale} again
\end{itemize}

\begin{tcolorbox}[breakable,colback=teal!4,colframe=teal!35!black,title={Before you move on}]
What does sale mean? (\textquotedblleft{}Dirty\textquotedblright{}.) Say it once more.
\end{tcolorbox}
\section[jeune]{jeune — young}

\label{lesson:FR-C92-jeune}



\begin{quote}
Before the new one: say the French for soft; mild; sweet, then the French for dirty.
\end{quote}

\subsection*{You'll want to know: jeune}
\textbf{jeune} — \textquotedblleft{}young\textquotedblright{}.

\begin{cousinweb}
\textbf{jeune} $\leftarrow$ Latin \emph{iuvenis}, \textquotedblleft{}young\textquotedblright{} (EtymDB 2.1 records the Latin word) — inherited.
\end{cousinweb}

\begin{grammarlens}[title={What you've built}]
One more word for this chapter.
\end{grammarlens}

\subsection*{Your turn}
\begin{itemize}
\raggedright
  \item \emph{Say it:} \emph{jeune}
  \item \emph{Say it:} \emph{jeune}, once more
  \item \emph{Say it:} say \emph{jeune}
  \item \emph{Recall:} say the French for soft; mild; sweet, then the French for dirty, then say \emph{jeune} again
\end{itemize}

\begin{tcolorbox}[breakable,colback=teal!4,colframe=teal!35!black,title={Before you move on}]
What does jeune mean? (\textquotedblleft{}Young\textquotedblright{}.) Say it once more.
\end{tcolorbox}
\section[gentil]{gentil — kind, nice}

\label{lesson:FR-C92-gentil}



\begin{quote}
Before the new one: say the French for dirty, then the French for young.
\end{quote}

\subsection*{You'll want to know: gentil}
\textbf{gentil} — \textquotedblleft{}kind, nice\textquotedblright{}.

\begin{cousinweb}
\textbf{gentil} $\leftarrow$ Latin \emph{gentilis}, \textquotedblleft{}of the same clan\textquotedblright{}, from \emph{gens}, \textquotedblleft{}a clan\textquotedblright{} — inherited. (EtymDB 2.1) In Old French it meant \textquotedblleft{}high-born, noble\textquotedblright{}.
\end{cousinweb}

\begin{grammarlens}[title={What you've built}]
One more word for this chapter.
\end{grammarlens}

\subsection*{Your turn}
\begin{itemize}
\raggedright
  \item \emph{Say it:} \emph{gentil}
  \item \emph{Say it:} \emph{gentil}, once more
  \item \emph{Say it:} say \emph{gentil}
  \item \emph{Recall:} say the French for dirty, then the French for young, then say \emph{gentil} again
\end{itemize}

\begin{tcolorbox}[breakable,colback=teal!4,colframe=teal!35!black,title={Before you move on}]
What does gentil mean? (\textquotedblleft{}Kind, nice\textquotedblright{}.) Say it once more.
\end{tcolorbox}
\section[étroit]{étroit — narrow}

\label{lesson:FR-C92-etroit}



\begin{quote}
Before the new one: say the French for young, then the French for kind.
\end{quote}

\subsection*{You'll want to know: étroit}
\textbf{étroit} — \textquotedblleft{}narrow\textquotedblright{}.

\begin{cousinweb}
\textbf{étroit} $\leftarrow$ Old French \emph{estroit} $\leftarrow$ Latin \emph{strictus}, \textquotedblleft{}drawn tight\textquotedblright{} (EtymDB 2.1) — inherited: \emph{-ct-} became \emph{-it}.
\end{cousinweb}

\begin{grammarlens}[title={What you've built}]
One more word for this chapter.
\end{grammarlens}

\subsection*{Your turn}
\begin{itemize}
\raggedright
  \item \emph{Say it:} \emph{étroit}
  \item \emph{Say it:} \emph{étroit}, once more
  \item \emph{Say it:} say \emph{étroit}
  \item \emph{Recall:} say the French for young, then the French for kind, then say \emph{étroit} again
\end{itemize}

\begin{tcolorbox}[breakable,colback=teal!4,colframe=teal!35!black,title={Before you move on}]
What does étroit mean? (\textquotedblleft{}Narrow\textquotedblright{}.) Say it once more.
\end{tcolorbox}
